\section{Implementation Details}
\label{sec:appendix_implementation}

\vspace{-0.25em}
\noindent
This appendix reports the \textbf{exact training stack}, \textbf{hyperparameters}, and \textbf{compute footprint} used to reproduce our results.
We emphasize two principles that guided the implementation:
\textbf{(i) comparability}---all methods share the same backbone, tokenization, LoRA target modules, and maximum context length; and
\textbf{(ii) auditability}---all reported metrics are deterministic, and all optimization settings are explicitly enumerated in Tables~\ref{tab:hardware}--\ref{tab:compute}.

% -----------------------------------------------------------------------------
\subsection{Hardware and Software}
% -----------------------------------------------------------------------------

\vspace{-0.25em}
\noindent
All experiments were run on NVIDIA A100 GPUs. Offline alignment methods (\textsc{DPO}, \textsc{CITA}) fit comfortably on \textbf{40GB} GPUs because they optimize on fixed preference pairs without on-policy generation.
In contrast, online methods (\textsc{PPO}, \textsc{GRPO}) required \textbf{80GB} GPUs due to the additional memory pressure from (i) generating multiple rollouts, (ii) caching token-level log-probabilities for policy-gradient updates, and (iii) maintaining reward/evaluator states during optimization.
We implement all methods in PyTorch using the HuggingFace Transformers ecosystem, and use TRL for standardized policy-optimization components (rollout, logprob tracking, KL control) to reduce implementation variance.

\begin{table*}[ht!]
\centering
\small
\renewcommand{\arraystretch}{1.2}
\begin{tabular}{@{}ll@{}}
\toprule
\textbf{Component} & \textbf{Specification} \\
\midrule
GPU & NVIDIA A100-40GB \\ \hline
Framework & PyTorch 2.1 + Transformers 4.35 \\ \hline
Training Library & TRL~\cite{tunstall2023zephyr} \\ \hline
Optimization & Optuna~\cite{akiba2019optuna} 3.4 with TPE sampler \\ \hline
Mixed Precision & bfloat16 \\ \hline
Gradient Checkpointing & Enabled \\
\bottomrule
\end{tabular}
\caption{\textbf{Hardware and software configuration.} We use TRL for consistent implementations of policy optimization and preference learning. Optuna with TPE is used for systematic hyperparameter search.}
\label{tab:hardware}
\end{table*}

\vspace{-0.35em}
\noindent\textbf{Implementation convention (shared across methods).}
We standardize (i) tokenizer and padding behavior, (ii) maximum sequence length (2048), and (iii) LoRA injection points to ensure differences arise from the \textbf{learning objective} rather than engineering artifacts.
All methods branch from the same SFT initialization; \textsc{DPO}/\textsc{CITA} then apply offline preference optimization, while \textsc{PPO}/\textsc{GRPO} run on-policy updates.

% -----------------------------------------------------------------------------
\subsection{Training Hyperparameters (Full)}
% -----------------------------------------------------------------------------

\vspace{-0.25em}
\noindent
Table~\ref{tab:full_hyperparams} reports the complete training hyperparameters used for each method.
We highlight three practical details that matter for reproduction:
\textbf{(i) effective batch parity}---we match effective batch sizes where feasible via gradient accumulation;
\textbf{(ii) identical context budgets}---max sequence length is held fixed across all runs;
\textbf{(iii) stable LoRA capacity}---LoRA rank/alpha and target modules are constant, so optimization difficulty is primarily dictated by the objective (offline contrast vs.\ online RL).

\vspace{-0.25em}
\noindent
\textbf{Why does \textsc{CITA} use a lower learning rate?}
\textsc{CITA\_Instruct} sequences are longer (due to explicit instruction conditioning and paired preference formatting), which increases gradient magnitude and variance.
We therefore select a lower learning rate for \textsc{CITA} than \textsc{DPO} to prevent overshooting and preserve a stable instruction-conditioned policy family.

\begin{table*}[ht!]
\centering
\scriptsize
\setlength{\tabcolsep}{2pt}
\renewcommand{\arraystretch}{1.8}
\begin{tabular}{@{}lcccc@{}}
\toprule
\textbf{Parameter} & \textbf{PPO} & \textbf{GRPO} & \textbf{DPO} & \textbf{CITA} \\
\midrule
Epochs & 1 & 1 & 1 & 1 \\ \hline
Batch Size & 16 & 12 & 1 & 1 \\ \hline
Grad. Accum. & 4 & 2 & 8 & 8 \\ \hline
Eff. Batch & 64 & 24 & 8 & 8 \\ \hline
Max Seq. Len. & 2048 & 2048 & 2048 & 2048 \\ \hline
Learning Rate & 1e-5 & 5e-6 & 1e-5 & 5.4e-6 \\ \hline
LR Scheduler & --- & Cosine & Cosine & Cosine \\ \hline
Warmup & --- & 10\% & 100 steps & 10\% \\ \hline
Weight Decay & --- & 0.1 & 0.01 & 0.011 \\ \hline
$\beta$ (temp.) & --- & --- & 0.1 & 0.107 \\ \hline
KL Coef. & 0.1 & --- & --- & 2.3e-4 \\ \hline
Max Grad Norm & 1.0 & 0.1 & --- & 1.0 \\ \hline
LoRA Rank & 16 & 16 & 16 & 16 \\ \hline
LoRA Alpha & 16 & 16 & 16 & 16 \\ \hline
Target Modules & \multicolumn{4}{c}{q, k, v, o, gate, up, down\_proj} \\
\bottomrule
\end{tabular}
\caption{\textbf{Complete hyperparameters for all methods.} PPO uses \texttt{mini\_batch=4}, \texttt{ppo\_epochs=4}, \texttt{cliprange=0.2}. GRPO uses \texttt{num\_generations=6}, \texttt{max\_completion=512}. LoRA configuration is fixed across methods to isolate objective-level effects.}
\label{tab:full_hyperparams}
\end{table*}

\vspace{-0.4em}
\noindent\textbf{Optuna and TPE (what it does here).}
We use Optuna’s \textbf{Tree-structured Parzen Estimator (TPE)} sampler to perform Bayesian-style hyperparameter search.
TPE fits separate density models over configurations that perform well vs.\ poorly and proposes new trials by maximizing an expected improvement criterion.
In our setting, we run 13 trials per method family and tune learning rate, weight decay, KL coefficient (when applicable), and $\beta$ for contrastive objectives.

% -----------------------------------------------------------------------------
\subsection{Compute Requirements}
% -----------------------------------------------------------------------------

\vspace{-0.25em}
\noindent
Table~\ref{tab:compute} summarizes the compute footprint.
The key comparison is \textbf{offline vs.\ online} training:
\textsc{DPO}/\textsc{CITA} operate on fixed preference pairs and therefore scale primarily with forward/backward passes,
whereas \textsc{PPO}/\textsc{GRPO} must \textbf{generate} completions during training, compute rewards, and store rollout statistics, which increases memory and wall-clock time.
Accordingly, \textsc{PPO}/\textsc{GRPO} require \textbf{A100-80GB} to avoid aggressive truncation or reduced rollout counts.

\begin{table*}[ht!]
\centering
\scriptsize
\renewcommand{\arraystretch}{1.8}
\begin{tabular}{@{}lcccc@{}}
\toprule
\textbf{Metric} & \textbf{PPO} & \textbf{GRPO} & \textbf{DPO} & \textbf{CITA} \\
\midrule
GPU Used & 80GB & 80GB & 40GB & 40GB \\ \hline
Training Time & 17h & 12h & 103min & 120min \\ \hline
GPU Memory & 72GB & 68GB & 39GB & 39GB \\ \hline
Checkpoint Size & 1.2GB & 1.2GB & 1.2GB & 1.2GB \\
\bottomrule
\end{tabular}
\caption{\textbf{Compute requirements per method.} \textsc{DPO}/\textsc{CITA} are offline and fit on A100-40GB; \textsc{PPO}/\textsc{GRPO} require A100-80GB primarily due to on-policy generation, rollout storage, and reward computation overhead.}
\label{tab:compute}
\end{table*}

\vspace{-0.35em}
\noindent\textbf{Interpretation of the footprint.}
The wall-clock gap is expected: online RL methods amortize additional cost per update by producing (and scoring) rollouts, while offline methods reuse a fixed dataset.
We therefore treat \textsc{PPO}/\textsc{GRPO} as \textbf{online reference points} for switchability rather than direct compute-matched baselines.

% -----------------------------------------------------------------------------
\subsection{Implementation Notes (Reproducibility-critical)}
% -----------------------------------------------------------------------------

\vspace{-0.25em}
\noindent
The following details are the most common sources of silent divergence in reproduction; we state them explicitly.

\begin{itemize}[leftmargin=*,itemsep=2.5pt]
    \item \textbf{Context construction (contract-first formatting).}
    We concatenate the alignment instruction $I$ and user request $X$ into a single context prefix, with consistent separators and role tags.
    This ensures the model treats $I$ as a \textbf{policy-level contract} rather than a stylistic hint.

    \item \textbf{Paired preference packing for \textsc{DPO}/\textsc{CITA}.}
    For each $(I,X)$ we compute $\log\pi_\theta(Y^+\mid I,X)$ and $\log\pi_\theta(Y^-\mid I,X)$ under the \emph{same} prefix.
    We verify that both completions fit within the maximum sequence length to avoid implicit truncation bias that can flip preference ordering.

    \item \textbf{Sequence-length discipline.}
    We keep \textbf{Max Seq.\ Len.\ = 2048} fixed across all methods.
    When $(I,X)$ is long, we truncate from the \emph{left} only after preserving the full instruction header, ensuring the behavioral contract is never partially dropped.

    \item \textbf{Hard negatives (when available).}
    When preference data provides multiple non-preferred candidates, we prioritize \textbf{hard negatives} (high-likelihood but wrong-under-contract) to sharpen $\Delta_\theta$ learning and reduce reliance on trivial stylistic cues.

    \item \textbf{KL control and stability.}
    For RL methods, a KL coefficient is used in the policy optimization loop to prevent policy explosion.
    For \textsc{CITA}, the KL anchor is used to keep instruction-conditioned regimes in a shared neighborhood.
    In preliminary sweeps, we found that overly large KL can suppress switching, while overly small KL increases regime interference; the reported value is selected by Optuna/TPE search.

    \item \textbf{Mixed precision + checkpointing.}
    We use \textbf{bfloat16} for stability and enable gradient checkpointing to fit long contexts without reducing batch size.
    All training runs are monitored for NaNs and exploding norms; when triggered, we lower LR and/or tighten gradient clipping.

    \item \textbf{LoRA injection (fixed capacity).}
    We apply LoRA with rank 16 to the same target modules across all methods (\texttt{q,k,v,o,gate,up,down\_proj}).
    This ensures gains in switchability are not due to differential parameter capacity.
\end{itemize}

\vspace{-0.25em}
\noindent
\textbf{Minimal reproduction recipe.}
To reproduce \textsc{CITA} results, it is sufficient to (i) start from the same SFT initialization, (ii) use the exact formatting for $(I,X)$, (iii) apply the Table~\ref{tab:full_hyperparams} settings with LoRA rank 16 on the specified modules, and (iv) evaluate with the deterministic metric suite reported in the main paper.




